\chapter{文件、设备与 I/O}

\section{原理}

文件描述符是用户态访问内核对象的句柄。它可以指向普通文件、管道、socket、设备或匿名内核对象。进程拿到的是一个小整数，但内核维护的是对象、offset、权限、引用计数和关闭状态。

先看一段最普通的代码：

\begin{lstlisting}
int fd = open("log.txt", O_CREAT | O_APPEND | O_WRONLY, 0644);
write(fd, "hello\n", 6);
close(fd);
\end{lstlisting}

初学时容易把它理解成三步：打开文件、把 6 个字节写到磁盘、关闭文件。真实路径更长：路径字符串要解析成目录项和 inode；\code{open} 要创建 open file description 并放入进程 fd 表；\code{write} 要验证 fd 权限、复制用户 buffer、更新文件 offset、修改 page cache 或提交设备请求；\code{close} 只是释放一个引用，不保证已经持久落盘。

把错误版本写出来更清楚：

\begin{lstlisting}
write(fd, buf, len):
  inode = fd_to_inode(fd)
  disk.write(inode.block, buf, len)
  return len
\end{lstlisting}

这个版本至少错在五处。第一，fd 不直接指向 inode，中间还有 open file description，它保存 offset 和 status flags。第二，\code{O\_APPEND} 必须在内核里原子决定写入位置，不能让用户态先 \code{lseek}。第三，普通文件写通常先进 page cache，\code{write} 返回不等于磁盘完成。第四，写入可能短写或被信号打断。第五，错误可能异步出现，例如后台 writeback 失败，之后 \code{fsync} 才报告。

文件 I/O 的第一层对象关系是：路径字符串经过路径解析找到目录项和 inode；\code{open} 产生 open file description；进程 fd 表保存一个小整数到 open file description 的引用。\code{dup} 产生新的 fd，但通常仍指向同一个 open file description，因此共享 offset 和状态。再次 \code{open} 同一路径通常产生新的 open file description，因此 offset 分离。

普通文件读写经常经过 page cache。读命中 page cache 时不碰设备；读未命中时可能提交块设备请求并等待完成。写通常先把用户字节复制进 page cache，把页标记为 dirty，再由后台 writeback 或 \code{fsync} 推进到设备。\code{write} 返回不等于数据安全落盘，这一点必须反复强调。

一次 buffered write 可以按这条线理解：

\begin{lstlisting}
sys_write(fd, user_buf, len):
  file = fget(fd)
  check file allows write
  decide offset, handling O_APPEND under lock
  copy bytes from user into page cache pages
  mark pages dirty
  update file size if needed
  return copied bytes

later:
  writeback finds dirty pages
  filesystem maps file offsets to blocks
  block layer submits requests
  device completes DMA
  fsync reports delayed writeback errors
\end{lstlisting}

这条线解释了三个常见困惑。第一，\code{write} 很快返回，是因为它可能只完成了“复制到内核缓存并标脏”。第二，系统掉电时，dirty page 还没写到持久介质就会丢。第三，\code{fsync} 的意义不是“再写一遍”，而是把先前的脏数据和必要元数据推进到更强的持久化边界，并把异步错误交还给调用者。

设备可以看作有特殊行为的文件对象。终端、磁盘、网卡、timerfd、eventfd、epoll、pipe、socket 都可以放进 fd/event 模型，但它们的 read/write 语义不同：有的返回字节流，有的返回结构化事件，有的可能阻塞，有的可能短读短写，有的只表示 readiness。

阻塞、非阻塞和异步是三种控制流模型。阻塞 I/O 让当前线程睡眠直到有结果；非阻塞 I/O 在不能推进时返回 \code{EAGAIN}，调用者等待 readiness 后重试；异步 I/O 提交请求，completion 稍后返回结果。三者没有绝对高低，关键是请求身份、buffer 生命周期、短读短写、错误传播和背压。

背压是 I/O 系统的安全阀。无界队列能让短 benchmark 看起来很快，因为生产者从不等待；一旦设备变慢，内存会持续增长，尾延迟扩大，最终进程可能被杀。可靠系统要同时限制 item 数、字节数和 in-flight 请求数，并在关闭时 drain 或 cancel。

\begin{keyidea}
fd 是入口，file 是打开实例，inode 是文件身份，page cache 是缓存和脏页账本，block request 是设备请求。文件 I/O 的连贯理解来自把这几层一条线串起来，而不是把 \code{read/write} 当成直接磁盘函数。
\end{keyidea}

\section{实践}

纸上实践一：画 fd 引用图。

\begin{lstlisting}
process fd table:
  fd 3 -> open file description A -> inode X
  fd 4 -> open file description A -> inode X   # dup
  fd 5 -> open file description B -> inode X   # open same path again
\end{lstlisting}

这张图能解释 offset 行为：fd 3 和 fd 4 共享 offset，fd 5 有独立 offset。并行文件读取如果使用 fd 3 和 fd 4，就可能因为共享 offset 得到不可预期的分片；使用 \code{pread} 或独立 open file description 更清楚。

纸上实践二：写出短写循环。任何 \code{write} 都可能只写入一部分字节，尤其是 pipe、socket、非阻塞 fd 或被信号打断时。正确代码必须记录已完成字节数，继续写剩余部分，遇到 \code{EINTR} 重试，遇到 \code{EAGAIN} 进入等待，而不是把短写当成完整成功。

纸上实践三：区分 readiness 和 message。\code{epoll} 告诉你 fd 可能可读或可写，不告诉你协议帧已经完整。网络程序必须有用户态 receive buffer、解析状态机和背压策略。把 fd readiness 当成完整请求，是事件驱动程序常见错误。

\begin{lstlisting}
PreparedBuffer
  -> Submitted
  -> InFlight
  -> Completed(bytes, error)
  -> ConsumedByProtocol
  -> BufferReleased
\end{lstlisting}

异步 I/O 中，buffer 在 completion 前不能被上游复用；completion 失败不能被提交路径吞掉；cancel 后也可能收到迟到 completion，所以请求身份必须稳定。

\section{算法与实现分析}

\subsection{fd 分配表}

进程的 fd table 可以先看作数组：下标是 fd，元素是指向 open file description 的引用。open 要找最小空闲 fd；close 清空槽位并减少引用计数；dup 找新槽位并增加引用计数。

\begin{lstlisting}
alloc_fd(file):
  for i in 0..max_fd:
    if table[i] == null:
      table[i] = file
      file.refcount += 1
      return i
  return -EMFILE

close(fd):
  file = table[fd]
  if file == null: return -EBADF
  table[fd] = null
  file.refcount -= 1
  if file.refcount == 0:
    file.ops.release(file)
\end{lstlisting}

线性扫描是 \code{O(n)}，但简单。真实系统会用 bitmap 或 hint 记录下一个可能空闲位置，把常见分配降到接近 \code{O(1)}。无论怎么优化，不变量都是：fd 只是进程局部索引，open file description 才保存 offset、flags 和对象引用。

\subsection{路径解析}

路径解析把字符串变成目录项和 inode。绝对路径从根目录开始，相对路径从当前工作目录开始；每个路径分量都要查目录；遇到符号链接可能要跳转；权限和 mount namespace 会改变可见结果。

\begin{lstlisting}
resolve(path):
  node = path.starts_with("/") ? root : current.cwd
  for component in split(path, "/"):
    if component == "." continue
    if component == "..": node = node.parent
    else:
      node = lookup_child(node, component)
      if node is symlink:
        node = resolve_symlink(node)
  return node
\end{lstlisting}

路径解析的复杂度大致和路径分量数相关，每个目录查找又取决于文件系统目录结构。缓存 dentry 可以避免每次都从磁盘读目录。安全实现还要限制符号链接递归深度，避免循环。

\subsection{page cache 读写}

普通文件读路径先查 page cache。若缓存命中，从内核页复制到用户 buffer；若未命中，分配 page cache 页，提交块 I/O，等待完成，再复制。写路径通常复制用户数据到 page cache，标记 dirty，并在之后由 writeback 或 fsync 推进设备写。

\begin{lstlisting}
buffered_read(file, offset, len):
  while len > 0:
    page_index = offset / PAGE_SIZE
    page = page_cache_lookup(file, page_index)
    if page == null:
      page = read_page_from_storage(file, page_index)
    n = copy_page_to_user(page, offset % PAGE_SIZE, len)
    offset += n
    len -= n
\end{lstlisting}

这个算法解释了冷读和热读差异。热读可能完全不碰设备；冷读可能等待存储。benchmark 若不说明缓存状态，就无法比较。

把一次冷读走完整。进程调用 \code{read(fd, buf, 6000)}，当前文件 offset 是 3000，页大小是 4096。第一次要读的是文件第 0 页的后 1096 字节，第二次要读第 1 页的前 4096 字节，最后还要读第 2 页的 808 字节。若这三页都不在 page cache，内核要为每页建立 page cache 条目，发起块 I/O，把线程挂到等待队列。设备 completion 到来后，页状态从 locked/loading 变成 uptodate，等待线程被唤醒，再把页内容复制到用户 buffer。

\begin{longtable}{p{0.16\linewidth}p{0.38\linewidth}p{0.34\linewidth}}
\toprule
阶段 & page cache 状态 & 线程看到什么 \\
\midrule
查第 0 页 & miss，分配 cache 页并标记 locked & 不能直接返回，需要等待 I/O \\
提交 I/O & 块层生成 request，设备 DMA 填页 & 线程 sleeping，不占 CPU \\
completion & 页标记 uptodate，清 locked，唤醒等待者 & 线程变 runnable，不一定立即运行 \\
复制用户 buffer & 从 page cache 页复制对应片段 & 可能再次遇到用户页 fault \\
推进 offset & 成功复制多少，offset 推进多少 & 短读时不能假设请求全完成 \\
\bottomrule
\end{longtable}

热读路径不同：若页已经 uptodate，内核只需要拿到页引用并复制到用户 buffer，不发设备 I/O。这里的“快”来自 page cache 保存了文件内容的内存副本，而不是文件系统跳过权限和边界检查。读性能测试若第一轮冷读、第二轮热读，却把两者当同一种读，就会得出错误结论。

\subsection{设备队列和完成}

块设备、网卡和异步 I/O 都可以抽象成提交队列和完成队列。提交时生成 request，放入设备或驱动队列；完成时中断或轮询拿到结果，唤醒等待任务或投递 completion。

\begin{lstlisting}
submit(req):
  req.id = next_id()
  req.state = IN_FLIGHT
  device_queue.push(req)
  kick_device()

complete(req_id, status, bytes):
  req = lookup(req_id)
  req.status = status
  req.bytes = bytes
  req.state = COMPLETED
  wake(req.waiter)
\end{lstlisting}

请求身份非常重要。异步完成可能乱序返回；取消后的请求也可能迟到完成；buffer 所有权必须跟着请求身份走。没有 request id，就无法把 completion 正确对应到提交者。

\subsection{fd、file 和 inode 的并发引用}

fd 是进程 fd 表下标，不是对象身份。内核执行 \code{read(fd)} 时，先把 fd 转成 \code{struct file} 引用；只要引用拿到，即使另一个线程随后 \code{close(fd)}，当前 I/O 也能在 file 对象上完成。close 删除的是 fd 表槽位，并减少 file 引用；它不能把正在执行的 I/O 用到的 file 直接释放。

\begin{lstlisting}
sys_read(fd, buf, len):
  file = fget(fd)
  if file == null:
    return -EBADF
  ret = vfs_read(file, buf, len)
  fput(file)
  return ret

sys_close(fd):
  file = files_table_remove(fd)
  if file == null:
    return -EBADF
  fput(file)
  return 0
\end{lstlisting}

这个规则解释 fd reuse 风险：线程 A 保存整数 fd=5，线程 B close(5)，线程 C open 得到新的 fd=5。A 再用整数 5 时，访问的可能已经是另一个对象。内核靠 file 引用保护内部生命周期，用户态也要避免把裸 fd 长期跨线程传递而没有所有权约定。

\subsection{open flags 的作用域}

flag 不是都绑定在同一层。fd flag 绑定 fd 表项；file status flag 绑定 open file description。dup 后两个 fd 指向同一个 open file description，所以共享 file status flag 和 offset；但 close-on-exec 是 fd flag，每个 fd 可以不同。

\begin{longtable}{p{0.20\linewidth}p{0.34\linewidth}p{0.34\linewidth}}
\toprule
标志 & 作用范围 & 边界 \\
\midrule
\code{O\_CLOEXEC} & fd 表项 & exec 提交时关闭该 fd，不影响其他 fd 的标志 \\
\code{O\_NONBLOCK} & open file description & dup 后共享，影响 read、write、accept、connect 的等待行为 \\
\code{O\_APPEND} & open file description & 每次写前在内核内原子移动到文件尾 \\
\code{O\_DIRECT} & file I/O 路径 & 对齐、页 pin、缓存一致性和短 I/O 更复杂 \\
\code{O\_SYNC/O\_DSYNC} & 写入持久化语义 & 返回前要推进数据或元数据到更强边界 \\
\bottomrule
\end{longtable}

\code{O\_APPEND} 不是用户态先 \code{lseek} 到末尾再 \code{write}。后者在并发下会竞争：两个线程都看到同一个旧文件尾，然后互相覆盖。append 语义要求内核在持有文件系统相关锁的状态下决定写入位置。

\subsection{短读短写不是异常}

普通文件顺序读常尽量返回请求长度，但很多 fd 都允许短读短写。pipe 可能只剩一部分数据；socket 是字节流，接收缓冲里有什么就先给什么；非阻塞 fd 无法继续推进时返回 \code{EAGAIN}；信号可能让系统调用返回 \code{EINTR}；存储错误可能让已经完成的部分和错误同时出现。

\begin{lstlisting}
write_all(fd, buf, len):
  done = 0
  while done < len:
    n = write(fd, buf + done, len - done)
    if n > 0:
      done += n
      continue
    if errno == EINTR:
      continue
    if errno == EAGAIN:
      wait_writable(fd)
      continue
    return error
  return done
\end{lstlisting}

读路径也需要协议层状态机。对 TCP，\code{recv} 返回 20 字节不表示一个业务消息完整；对 UDP，一次 \code{recvmsg} 对应一个 datagram，buffer 太小会截断；对终端，行规程可能改变返回时机。OS 提供 fd 语义，应用协议必须自己定义消息边界。

短写也要用具体事故理解。假设应用要把 1MiB 日志写到非阻塞 socket，第一次 \code{write} 返回 64KiB，第二次返回 \code{EAGAIN}。如果应用把第一次返回当作“整条日志已经发送”，就会丢掉后 960KiB；如果它在 \code{EAGAIN} 时忙循环重试，会把 CPU 打满；如果它没有保存剩余 buffer 和当前偏移，等 socket 可写时也不知道从哪里继续。正确实现需要一个输出队列，记录每条待发送数据、已发送 offset 和连接关闭时的清理策略。

\begin{lstlisting}
connection_on_writable(conn):
  while !conn.outq.empty():
    item = conn.outq.front()
    n = write(conn.fd, item.buf + item.off, item.len - item.off)
    if n > 0:
      item.off += n
      if item.off == item.len:
        conn.outq.pop_front()
      continue
    if errno == EAGAIN:
      arm_epollout(conn)
      return
    if errno == EINTR:
      continue
    fail_connection(conn, errno)
    return
\end{lstlisting}

这段代码把“短写不是异常”落实成状态机：数据没写完时必须保留；\code{EAGAIN} 表示当前 fd 暂不可推进；真正错误才关闭连接或上报。文件 fd、pipe、socket、终端的短 I/O 原因不同，但调用者都不能假设一次系统调用等于业务消息完整提交。

\subsection{buffered I/O 和 direct I/O 的不同风险}

buffered I/O 经过 page cache，适合普通应用和共享缓存；direct I/O 尽量绕过 page cache，适合数据库一类自己管理缓存和写入顺序的程序。direct I/O 不是“更高级的 read/write”，它把对齐、页 pin、设备队列、缓存一致性和错误处理暴露得更多。

\begin{longtable}{p{0.22\linewidth}p{0.34\linewidth}p{0.32\linewidth}}
\toprule
路径 & 优点 & 代价 \\
\midrule
buffered read & 命中 page cache 很快，可共享缓存 & 多一次复制，缓存污染可能影响其他任务 \\
buffered write & 快速标脏，后台批量写回 & write 返回不代表持久化 \\
direct read/write & 减少 page cache 干扰，适合自管缓存 & 对齐、pin 页、短 I/O 和并发一致性复杂 \\
mmap & load/store 直接访问映射页 & 缺页、SIGBUS、dirty/writeback 语义更隐蔽 \\
\bottomrule
\end{longtable}

direct I/O 和 buffered I/O 混用尤其危险。一个路径绕过 page cache，另一个路径读写 page cache，若文件系统没有正确失效或同步缓存，应用可能看到旧数据。真实系统会在文件系统层处理这类一致性，但应用设计仍应尽量避免无计划混用。

\subsection{epoll item 绑定的是 file}

\code{epoll\_ctl} 把一个 fd 对应的 file 和 interest mask 加进 epoll 实例。之后整数 fd 关闭并复用，不代表旧 epoll item 自动变成新文件的监听项。事件返回里带着用户当初登记的 fd 数字，但内核里的等待关系绑定的是旧 file。

\begin{lstlisting}
epoll_ctl(epfd, ADD, fd):
  file = fget(fd)
  epitem.file = file
  epitem.fd_number_for_user = fd
  attach_poll_callback(file, epitem)

event delivery:
  file state changes
  callback enqueues epitem
  epoll_wait returns stored fd_number and event mask
\end{lstlisting}

事件循环通常要给每个连接对象加 generation 或指针身份，而不是只根据 fd 整数找状态。close、dup、fork、exec、线程并发都可能让 fd 数字的直觉失效。

\subsection{io\_uring 的提交、完成和取消}

io\_uring 把“提交请求”和“收完成结果”拆成两个环。SQE 中的 \code{user\_data} 是应用识别 completion 的关键；CQE 中的 \code{res} 是最终字节数或负 errno。提交成功只表示内核接收了请求，不表示 I/O 完成。

\begin{lstlisting}
submit:
  fill SQE with opcode, fd, buffer, length, user_data
  publish SQ tail
  enter kernel or rely on SQPOLL

complete:
  read CQE
  id = cqe.user_data
  result = cqe.res
  release or reuse buffer only now
\end{lstlisting}

取消请求要按竞态理解：取消到达时，请求可能还在队列里，可能已经派发给设备，可能刚刚完成。应用必须能接受“取消成功”“取消失败因为已经完成”“先收到完成后收到取消结果”等情况。只要 completion 可能到来，buffer 和请求对象就不能提前复用。

更具体地说，\code{io\_uring} 取消不是“把那次 I/O 从世界上抹掉”。假设应用提交读请求 R，\code{user\_data=42}，buffer 指向一块复用池内存。100ms 后应用超时，提交 cancel(42)。此时可能有三种真实状态：R 还在内核队列，取消可以把它摘掉；R 已经交给块设备，设备可能无法撤销，只能等 completion；R 已经完成，CQE 还没被应用收走。应用若在提交 cancel 后立刻把 buffer 给另一个请求使用，就可能和迟到的 R completion 冲突。

\begin{longtable}{p{0.18\linewidth}p{0.38\linewidth}p{0.32\linewidth}}
\toprule
取消到达时 & 内核可能返回什么 & 应用必须怎样处理 \\
\midrule
请求未派发 & cancel CQE 表示取消成功，原请求不会再完成 & 可以释放 buffer，但仍要消费 cancel 结果 \\
请求已派发 & cancel 可能失败或标记取消中，原请求仍会完成 & buffer 继续保持有效，直到看到原请求 CQE \\
请求已完成 & cancel 返回未找到或已完成 & 根据原 CQE 的结果处理，不重复释放 \\
完成和取消乱序到达 & 先看到哪一个都可能 & 用 request 状态机合并结果，避免双重释放 \\
\bottomrule
\end{longtable}

所以异步 I/O 的基本规则是：提交记录、buffer 所有权和 completion 消费必须绑定。真正能释放或复用 buffer 的点，不是“我不想等了”，而是“所有可能引用这个 buffer 的请求状态都已经收敛”。这和驱动层 DMA request 的规则是同一件事，只是对象从设备 descriptor 换成了用户可见的 SQE/CQE。

\subsection{源码阅读入口}

fd 表和引用看 \filepath{fs/file.c}、\filepath{include/linux/fdtable.h}；open 和路径解析看 \filepath{fs/open.c}、\filepath{fs/namei.c}；read/write 主干看 \filepath{fs/read\_write.c}；page cache 看 \filepath{mm/filemap.c}；epoll 看 \filepath{fs/eventpoll.c}；io\_uring 看 \filepath{io\_uring/}。阅读时跟住一个 fd：它从整数变成 file，如何增加引用，如何进入具体 file operations，错误码在哪里返回，最后一个引用在哪里释放。

\section{测试}

测试覆盖 open/write/read、dup 后共享内容、close 一个 fd 后另一个 fd 仍有效、关闭所有引用后访问失败、未知 fd 返回诊断。还要测试 fd 分配单调或复用策略是否有明确规则。

\begin{itemize}
\tightlist
\item \code{dup} 后 offset 共享，独立 \code{open} 后 offset 分离。
\item \code{pread} 不改变 current offset。
\item pipe 空读阻塞或非阻塞返回 \code{EAGAIN}，写端关闭后读端得到 EOF。
\item 读端关闭后写 pipe 失败，不应静默丢数据。
\item socket 或 pipe 短写后，调用者继续写剩余字节。
\item readiness 事件只触发解析尝试，不被当成完整消息。
\item 异步请求完成前，buffer 所有权不能释放。
\item 队列达到字节上限时，生产者受到背压。
\end{itemize}

第一版必须先保证句柄生命周期可解释。持久化和崩溃恢复不要另起一条线，而要回到 fd、file、inode、page cache、block request 和目录项这些对象上验收。

\subsection{本章压缩进来的 VFS、文件系统与恢复}

VFS 给用户一个统一 fd 接口，具体文件系统负责目录项、inode、块映射、权限、日志和恢复。路径解析得到 dentry/inode，open 生成 file，read/write 进入 file operations，普通文件再落到 address\_space、page cache 和文件系统块映射。主线里先抓住对象边界，而不是先背某个文件系统格式。

\begin{longtable}{p{0.20\linewidth}p{0.38\linewidth}p{0.30\linewidth}}
\toprule
对象 & 保存什么 & 关键不变量 \\
\midrule
dentry & 名字到 inode 的缓存关系 & rename/unlink 后缓存不能让旧名字错误复活 \\
inode & 文件身份、权限、大小、块映射、时间戳 & 多个路径和 fd 可引用同一 inode \\
file & 一次打开实例、offset、status flags、ops & \code{dup} 共享，独立 \code{open} 分离 \\
address\_space & 文件页缓存和写回状态 & read/write/mmap 必须看到同一份 page cache 账本 \\
journal/log & 元数据或数据提交记录 & 崩溃恢复只能重放完整事务 \\
\bottomrule
\end{longtable}

文件系统恢复的核心不是“写得更勤”，而是“崩溃后能判断哪些写已经构成完整状态”。以创建文件并重命名为例，数据块、inode 大小、目录项、父目录元数据、日志 commit 和设备 flush 之间必须有顺序。若应用写临时文件、\code{fsync(tmp)}、\code{rename(tmp, final)}，但忘记 \code{fsync} 父目录，崩溃后目录项是否存在就依赖文件系统和挂载语义，不能当作通用保证。

\begin{lstlisting}
durable_replace(tmp, final):
  write tmp contents
  fsync(tmp_fd)
  rename(tmp, final)
  fsync(parent_dir_fd)
\end{lstlisting}

这段不是所有场景的万能配方，但它展示了恢复协议的写法：每一步都要说明崩溃发生后恢复程序能看到什么。若看到旧文件，旧版本必须仍完整；若看到新名字，新内容必须已经同步；若看到临时文件，恢复程序应能删除或继续提交。文件系统、块层和设备只提供若干提交边界，业务正确性还要靠应用自己的版本号、校验和、manifest 或事务日志。

\begin{rubricbox}
文件与 I/O 章的最终验收：跟住一个字节从用户 buffer 到 page cache、bio、request、设备 completion、writeback error、\code{fsync} 返回和恢复程序判断。若不能指出每层的状态对象和错误返回点，就不能声称理解了文件 I/O。
\end{rubricbox}
